\documentclass[10pt,twocolumn,letterpaper]{article}

\usepackage{graphicx}
\usepackage{amsmath,amssymb}
\usepackage{booktabs}
\usepackage{multirow}
\usepackage{hyperref}
\usepackage{xcolor}
\usepackage{subcaption}

% Convenience
\newcommand{\eg}{\emph{e.g.}}
\newcommand{\ie}{\emph{i.e.}}
\newcommand{\etal}{\emph{et al.}}

\title{Beyond Single Prompts: Comparing Vision Adapters, CoOp,\\
and Optimal Transport for Domain-Specific Image Classification}

\author{
  Mattia Lomuto (VR543327)
  \and Carlo Stast (VR543606)
  \and Marco Equisetto (VR535007)
  \\[4pt]
  \textit{Advanced Programming for AI}\\
  \textit{[University of Verona]}\\
  {\small\texttt{\{mattia.lomuto, carlo.stasi, marco.equisetto\}@studenti.univr.it}}
}

\begin{document}
\maketitle

% ============================================================================
% 1. ABSTRACT
% ============================================================================
\begin{abstract}

Large-scale vision-language models such as CLIP offer strong zero-shot transfer, yet their performance on specialised visual domains (like satellite imagery, textures, fine-grained flowers) still falls behind lightweight task-specific adaptation.  The question is not only \emph{which adapter is most accurate}, but also \emph{how much accuracy each trainable parameter buys}.

We present a controlled comparison of seven parameter-efficient adaptation strategies applied to a \emph{single frozen} CLIP ViT-B/32 backbone:

(i)~zero-shot prompting and prompt ensembling,
(ii)~linear probing,
(iii)~Context Optimization (CoOp) with learned text prompts,
(iv)~CLIP-Adapter with a vision-side bottleneck MLP,
(v)~a joint CoOp\,+\,Adapter model that adapts both modalities simultaneously,
(vi)~Tip-Adapter and its fine-tuned variant Tip-Adapter-F with a training-free
key--value cache, and
(vii)~Low-Rank Adaptation (LoRA) injected into the vision tower.

As an exploratory baseline, we also evaluate Sinkhorn-based Optimal Transport (OT) on patch-level tokens, which underperforms zero-shot on all benchmarks and provides instructive negative results.

% TODO: Update numbers after final verified run
All methods share the same frozen backbone, the same evaluation loop, and the same few-shot and full-shot protocol across three datasets (EuroSAT, DTD, Flowers102), making accuracy, parameter count, GPU memory, and wall-clock time directly comparable. Our main finding is that CoOp's 8\,192 learnable prompt tokens, which is 0.0095\,\% of the backbone, close the majority of the gap between zero-shot and full supervision, while Tip-Adapter-F achieves the highest overall accuracy
with a fine-tuned visual cache.

\smallskip\noindent\textbf{Code:}~\url{https://github.com/MarcoEquisetto/APAI_exam}
\end{abstract}


\section{Introduction}\label{sec:introduction}

\subsection{Problem Statement and Motivation}

Contrastive Language-Image Pre-training (CLIP)~\cite{radford2021learning} has demonstrated remarkable zero-shot transfer capabilities across various visual recognition tasks.  By aligning images and captions in a shared embedding space, CLIP enables classification via natural-language prompts without any task-specific training. Still, in specific domains
such as remote sensing, material textures and fine-grained botany, zero-shot accuracy remains far from satisfactory: a frozen ViT-B/32 achieves only 46.4\,\% on EuroSAT~\cite{helber2019eurosat}, 54.9\,\% on DTD~\cite{cimpoi2014describing}, and 71.5\,\% on Flowers102~\cite{nilsback2008automated}, well below the needs of any practical application.

The standard remedy is full fine-tuning, but updating all 86~million parameters of a ViT-B/32 requires large labelled datasets and risks catastrophic forgetting of the broad visual and linguistic knowledge that makes CLIP useful in the first place, without even considering the computational toll of such large scale updates. A rapidly growing family of 
\emph{parameter-efficient adaptation} methods addresses this by training a small number of new parameters while keeping the backbone frozen. These methods differ in \emph{where} they intervene and in \emph{how many parameters} they introduce, from zero (training-free caches) to several hundred thousand (LoRA inside the backbone).


\subsection{Challenges and Prior Work}

Several adaptation strategies have already been proposed for CLIP. \textbf{CoOp}~\cite{zhou2022learning} replaces the hand-written text template with $M$ continuous context vectors optimised by gradient descent.
\textbf{CLIP-Adapter}~\cite{gao2024clipadapter} appends a bottleneck MLP after the frozen visual encoder and connects adapted features with the originals via a residual connection.
\textbf{Tip-Adapter}~\cite{zhang2022tip} takes a training-free approach, constructing a key-value cache from few-shot support images and combining cache-based logits with zero-shot predictions. Its fine-tuned variant, \textbf{Tip-Adapter-F}, further optimises the cache keys by gradient descent.
\textbf{LoRA}~\cite{hu2022lora}, originally created for large language models, injects low-rank matrices into existing linear layers, enabling adaptation \emph{inside} the backbone with a tollerable parameter overhead.
Finally, \textbf{Optimal Transport}~\cite{cuturi2013sinkhorn} offers a principled framework for comparing distributions of tokens, using the \textbf{Sinkhorn} algorithm to compute a structured distance between patch-level image features and word-level text features.

Despite this solution-rich environment, it appears existing studies typically evaluate each method in isolation, with different backbones, different datasets, different shot counts, and different hyperparameter protocols, making fair comparisons
difficult.


\subsection{Limitations of Existing Comparisons}

\begin{enumerate}
  \item \textbf{Inconsistent backbones.} CoOp reports results on a ResNet-50 CLIP, while CLIP-Adapter uses ViT-B/16 and Tip-Adapter uses ViT-B/32... all under different pre-training regimes.
  
  \item \textbf{Uncontrolled shot protocols.} Methods are often compared at different shot counts or with different sampling seeds, confounding accuracy differences with data selection effects.
  
  \item \textbf{Missing efficiency metrics.} Papers report accuracy but rarely tabulate GPU memory, training time, or the number of trainable parameters on the same hardware, leaving users with no basis for cost-benefit analysis.
  
  \item \textbf{No negative results.}  Methods that underperform zero-shot CLIP (such as naïve Optimal Transport on token features) are typically unpublished, depriving of insights into failure modes.
\end{enumerate}


\subsection{Contributions}

We make the following contributions:

\begin{enumerate}
  \item \textbf{Unified benchmark.} We implement ten adaptation strategies within a single codebase, all sharing the same frozen ViT-B/32 backbone (LAION-2B weights via OpenCLIP), the same evaluation loop, the same few-shot sampling, and  the same epoch budget. This makes accuracy, parameter count, peak GPU memory, and wall-clock time directly comparable across every method.

  \item \textbf{Efficiency analysis.} We attempt to quantify the \emph{accuracy per trainable parameter} of each method across three domains (EuroSAT, DTD, Flowers102), revealing that the majority of the zero-shot-to-supervised gap can be closed with fewer than 10\,000 parameters, and that the marginal return per parameter decreases sharply beyond that
  point.

  \item \textbf{Joint adaptation.} We construct a joint CoOp\,+\,CLIP-Adapter model that adapts both text and vision
  modalities simultaneously, and document the software engineering required to compose two independently developed adapter
  modules (MRO conflicts, optimizer group splits, gradient flow traps).

  \item \textbf{Informative negative result.} We provide a thorough analysis of why Sinkhorn-based Optimal Transport on patch-level features fails to outperform zero-shot CLIP: patch tokens are not in the contrastive aligned space, uniform marginals over-weight background patches, and shared template words dominate the text-token distribution.

  \item \textbf{Open-source reproducibility.} The complete codebase, including all experiment scripts, plotting, and raw result files, is publicly available.
  
\end{enumerate}


\subsection{Paper Overview}

The remainder of this paper is organized as follows.
Section~\ref{sec:method} details our experimental framework, mathematical formulations, and training protocol.
Section~\ref{sec:results} presents comparative results and ablation studies.
Section~\ref{sec:conclusion} summarizes our findings, discusses limitations, and suggests directions for future work.
Section~\ref{sec:contributions} details the individual contributions of each team member.


\section{Method}\label{sec:method}

% ----- PLACEHOLDER: Pipeline Architecture Diagram -----
% TODO: INSERT PIPELINE ARCHITECTURE DIAGRAM HERE
% The diagram should show: frozen CLIP backbone (ViT-B/32) in the centre,
% with the text encoder on one side and the vision encoder on the other.
% Each adaptation method should be shown as a coloured module plugged into
% the appropriate encoder output.  Suggested tool: TikZ or draw.io.
%
% \begin{figure*}[t]
%   \centering
%   \includegraphics[width=\textwidth]{figures/pipeline.pdf}
%   \caption{Overview of the experimental pipeline.  All methods share a
%     single frozen CLIP ViT-B/32 backbone and differ only in the coloured
%     adapter modules.}
%   \label{fig:pipeline}
% \end{figure*}
% -------------------------------------------------------


\subsection{Frozen Backbone}\label{sec:backbone}

All experiments use a ViT-B/32 backbone loaded from OpenCLIP with LAION-2B weights (\texttt{laion2b\_s34b\_b79k}). The backbone comprises a 12-layer Vision Transformer (86\,M parameters) producing 512-dimensional embeddings for both images and text. All backbone parameters are frozen (\texttt{requires\_grad\,=\,False}); only parameters introduced by each
adaptation module are trainable.

A shared base class, \texttt{BaseCLIPWrapper}, encapsulates the frozen backbone and exposes two virtual methods, namely texttt{get\_image\_features()} and \texttt{get\_text\_features()}, that subclasses override to inject their adapter logic. A single \texttt{predict()} method, inherited by all adapters, computes cosine similarities between image and text
prototypes, ensuring that every method is evaluated through the same code path.


\subsection{Baseline Methods}\label{sec:baselines}

\paragraph{Zero-Shot CLIP.}
Classification is performed by encoding each class name with a prompt template (\eg, ``a satellite image of \{class\}'') and selecting the class whose text embedding has the highest cosine similarity to the image embedding. No parameters are trained.

\paragraph{Zero-Shot with Prompt Ensembling.}
Following Radford~\etal~\cite{radford2021learning}, we average the text prototypes produced by 7~templates that vary in viewpoint and framing (\eg, ``a satellite photo of \{\}'', ``an aerial view of \{\}''). This costs no additional parameters and typically adds 2-4 accuracy points.

\paragraph{Linear Probe.}
An $\ell_2$-regularised logistic regression (implemented via \texttt{scikit-learn}) is fitted on the frozen image features.  The number of trainable parameters is $C \times (D + 1)$, where $C$ is the number of classes and $D = 512$.


\subsection{CoOp: Learnable Text Prompts}\label{sec:coop}

Context Optimization (CoOp)~\cite{zhou2022learning} replaces the hand-written template with $M$ continuous context vectors in the text transformer's input embedding space:

\begin{equation}\label{eq:coop}
  \mathbf{p}_c = [\mathrm{SOS}]\;
  \underbrace{\mathbf{v}_1\;\mathbf{v}_2\;\cdots\;\mathbf{v}_M}_{\text{learnable}}
  \;[\mathrm{tok}(c)]\;[\mathrm{EOS}]
\end{equation}

where $[\mathrm{tok}(c)]$ denotes the frozen BPE embedding of class name $c$. In the \emph{Unified Context} variant a single $(\mathbf{v}_1, \ldots, \mathbf{v}_M)$ block is shared across all classes; in the \emph{Class-Specific Context} (CSC) variant each class receives its own block, multiplying the parameter count by $C$.

For ViT-B/32 with $M = 16$, the unified variant introduces $16 \times 512 = 8{,}192$ trainable parameters. The text transformer is run from the assembled embeddings directly, skipping the frozen token embedding table for the context positions, so gradients flow only through $\mathbf{v}_1, \ldots, \mathbf{v}_M$.  We train with SGD (momentum 0.9,
learning rate $2 \times 10^{-3}$) and a cosine-annealing schedule with a linear warmup epoch.


\subsection{CLIP-Adapter: Vision-Side Bottleneck MLP}\label{sec:adapter}

CLIP-Adapter~\cite{gao2024clipadapter} appends a trainable bottleneck MLP after the frozen vision encoder:

\begin{align}
  \mathbf{f}_{\mathrm{adapt}} &= W_{\mathrm{up}} \;\sigma\bigl(
    W_{\mathrm{down}}\,\mathbf{f}\bigr) \label{eq:adapter-mlp}\\
  \tilde{\mathbf{f}} &= \alpha\,\mathbf{f}_{\mathrm{adapt}}
    + (1-\alpha)\,\mathbf{f} \label{eq:adapter-blend}
\end{align}

where $\mathbf{f} \in \mathbb{R}^D$ is the frozen, L2-normalised image embedding, $W_{\mathrm{down}} \in \mathbb{R}^{(D/r) \times D}$ and $W_{\mathrm{up}} \in \mathbb{R}^{D \times (D/r)}$ form the bottleneck
with reduction ratio $r$, $\sigma$ is ReLU, and $\alpha \in [0, 1]$ is a residual blending factor.
The adapted features $\tilde{\mathbf{f}}$ are L2-normalised before computing similarity.

With $r = 4$ and $D = 512$ the adapter introduces $131{,}712$ trainable parameters. $W_{\mathrm{up}}$ is initialised near zero so that training begins close to the original CLIP representation. We train with AdamW (learning rate $10^{-3}$, cosine schedule).


\subsection{Joint CoOp + CLIP-Adapter}\label{sec:joint}

The two methods are complementary by construction: CoOp overrides \texttt{get\_text\_features()} and CLIP-Adapter overrides \texttt{get\_image\_features()}. We compose them in a single model that adapts both modalities simultaneously, for a total of $8{,}192 + 131{,}712 = 139{,}904$ parameters.

Naïve Python multiple inheritance fails because of a Method Resolution Order (MRO) conflict: during construction,
\texttt{CLIPAdapterModel.\_\_init\_\_()} calls \texttt{self.\_update\_text\_prototypes()}, which resolves to CoOp's text
path and reads \texttt{self.prompt\_learner}, which is a module that has not been created yet. We resolve this by constructing both modules explicitly in a controlled order and assigning the feature methods as function attributes.

Crucially, the two halves require different optimizer recipes. CoOp trains with SGD at $2 \times 10^{-3}$; the adapter trains with AdamW at $10^{-3}$.  A single learning rate applied to both degrades accuracy below either method alone. We therefore define per-module parameter groups under one AdamW optimizer.


\subsection{Tip-Adapter and Tip-Adapter-F}\label{sec:tip}

Tip-Adapter~\cite{zhang2022tip} constructs a training-free key-value cache $(\mathbf{F}_{\mathrm{train}}, \mathbf{L}_{\mathrm{train}})$ from the few-shot support set, where $\mathbf{F}_{\mathrm{train}} \in
\mathbb{R}^{N \times D}$ are frozen image features and $\mathbf{L}_{\mathrm{train}} \in \mathbb{R}^{N \times C}$ are one-hot label vectors. At test time:
\begin{equation}\label{eq:tip}
  \mathrm{logits} = 100\,\mathbf{f}_{\mathrm{test}}
    \mathbf{W}_{\mathrm{text}}^\top
  + \alpha\,\exp\!\bigl(-\beta(1 - \mathbf{f}_{\mathrm{test}}
    \mathbf{F}_{\mathrm{train}}^\top)\bigr)\,\mathbf{L}_{\mathrm{train}}
\end{equation}
where the first term is the zero-shot logit and the second is a cache-based correction weighted by visual similarity.

\textbf{Tip-Adapter-F} promotes the cache keys $\mathbf{F}_{\mathrm{train}}$ to learnable parameters and fine-tunes them with cross-entropy loss (AdamW, $\mathrm{lr} = 10^{-3}$, 20 epochs).  The number of trainable parameters is $N \times D = 16 \cdot C \cdot 512$, which scales with the number of classes.


\subsection{LoRA: Low-Rank Adaptation Inside the Backbone}\label{sec:lora}

Low-Rank Adaptation (LoRA)~\cite{hu2022lora} injects trainable low-rank matrices into existing linear layers.  For each target layer with weight $W_0 \in \mathbb{R}^{d_{\mathrm{out}} \times d_{\mathrm{in}}}$, LoRA adds:

\begin{equation}\label{eq:lora}
  W = W_0 + \frac{\alpha}{r}\,B\,A
\end{equation}

where $A \in \mathbb{R}^{r \times d_{\mathrm{in}}}$ and $B \in \mathbb{R}^{d_{\mathrm{out}} \times r}$ with rank $r \ll
\min(d_{\mathrm{in}}, d_{\mathrm{out}})$; $\alpha$ scales the contribution. $W_0$ remains frozen.

We apply LoRA to the MLP layers (\texttt{c\_fc} and \texttt{c\_proj}) of every transformer block in the vision tower with $r = 4$ and $\alpha = 1.0$, introducing $368{,}640$ trainable parameters. Unlike all other methods in this study, LoRA modifies the computation \emph{inside} the backbone; the pretrained weights still remain frozen, but this is a
deliberate out-of-brief comparison.  We train with AdamW at $\mathrm{lr} = 10^{-4}$ (10$\times$ lower than the adapter) with cosine scheduling.

\subsection{Optimal Transport (Exploratory)}\label{sec:ot}

As an exploratory baseline, we use Sinkhorn-based Optimal Transport~\cite{cuturi2013sinkhorn} to compute structured distances between the full set of image patch tokens and text word tokens. For an image with $P$ patch tokens $\{\mathbf{x}_i\}$ and a class with $T$ text tokens $\{\mathbf{y}_j\}$, we build a cost matrix $C_{ij} = 1 - \cos(\mathbf{x}_i, \mathbf{y}_j)$ and solve: 
\begin{equation}\label{eq:sinkhorn}
  d_{\varepsilon}(\mathbf{a}, \mathbf{b})
    = \min_{\gamma \in \Pi(\mathbf{a}, \mathbf{b})}
      \sum_{ij} \gamma_{ij}\,C_{ij}
      + \varepsilon\,\mathrm{KL}(\gamma \,\|\, \mathbf{a} \otimes \mathbf{b})
\end{equation}

with uniform marginals $\mathbf{a} = \frac{1}{P}\mathbf{1}$ and $\mathbf{b} = \frac{1}{T}\mathbf{1}$, regularisation $\varepsilon = 0.05$, and maximum 100 Sinkhorn iterations.  The predicted class is the one with the lowest transport cost.  No parameters are trained.


\subsection{Datasets}\label{sec:datasets}

We evaluate on three benchmarks spanning distinct visual domains:

\begin{itemize}
  \item \textbf{EuroSAT}~\cite{helber2019eurosat}: 27{,}000 Sentinel-2 satellite images, 10 land-use classes.  Template: ``a satellite image of \{class\}''.
  
  \item \textbf{DTD}~\cite{cimpoi2014describing}: 5{,}640 texture images, 47 classes.  Template: ``a photo of a \{class\} texture''.
  
  \item \textbf{Flowers102}~\cite{nilsback2008automated}: 8{,}189 flower images, 102 fine-grained categories.  Template: ``a photo of a \{class\}, a type of flower''.
\end{itemize}

All images are resized to $224 \times 224$ with bicubic interpolation and normalised with ImageNet statistics.  No data augmentation is applied.


\subsection{Training Protocol}\label{sec:protocol}

All trainable methods use 16-shot few-shot sampling with balanced class representation, a matched epoch budget of 10 epochs, batch size 64, and a cosine-annealing learning rate schedule with a 1-epoch linear warmup.
% TODO: Verify if your final run used full-shot instead of 16-shot.
% The unified_results.json config says shots=16.
Experiments are run on an NVIDIA GPU with PyTorch~2.5.1 (CUDA~12.1).


% ============================================================================
% 4. RESULTS
% ============================================================================
\section{Results}\label{sec:results}

% TODO: VERIFY ALL NUMBERS BELOW against your final run before submission.
% The numbers below come from unified_results.json (shots=16, epochs=10,
% seed=0).

\subsection{Main Comparison}\label{sec:main-results}

Table~\ref{tab:main} reports top-1 accuracy, trainable parameter count, peak GPU memory, and training time for all methods on the three datasets.

\begin{table*}[t]
\centering
\caption{Top-1 accuracy (\%), trainable parameters, peak GPU memory (MB), and training time (s) across three datasets.  All methods use a frozen ViT-B/32 backbone. Best accuracy in \textbf{bold}; second-best \underline{underlined}.}

\label{tab:main}
\small
\begin{tabular}{@{}lrrrrrrr@{}}
\toprule
\multirow{2}{*}{\textbf{Method}} &
\multirow{2}{*}{\textbf{Params}} &
\multicolumn{2}{c}{\textbf{EuroSAT}} &
\multicolumn{2}{c}{\textbf{DTD}} &
\multicolumn{2}{c}{\textbf{Flowers102}} \\
\cmidrule(lr){3-4}\cmidrule(lr){5-6}\cmidrule(lr){7-8}
& & Acc. & Mem. & Acc. & Mem. & Acc. & Mem. \\
\midrule
Zero-Shot           & 0       & 46.35 & 738  & 54.89 & 747  & 71.46 & 747  \\
ZS + Ensembling     & 0       & 50.61 & 738  & 57.34 & 747  & 71.93 & 747  \\
Linear Probe        & 5\,130  & 78.94 & 738  & 65.16 & 747  & 86.97 & 747  \\
Optimal Transport   & 0       & 20.31 & 738  & 23.30 & 747  & 19.06 & 749  \\
\midrule
CoOp                & 8\,192  & 75.87 & 1\,424  & 72.45 & 2\,138  & 88.71 & 3\,213  \\
CLIP-Adapter        & 131\,712 & 81.26 & 1\,347  & 74.31 & 1\,339  & 91.72 & 1\,420  \\
CoOp + Adapter      & 139\,904 & 81.59 & 1\,426  & 74.15 & 2\,140  & 91.49 & 3\,215  \\
\midrule
Tip-Adapter         & 0       & 74.41 & 1\,327  & 62.07 & 1\,329  & 81.17 & 1\,330  \\
Tip-Adapter-F       & 81\,920 & \textbf{86.26} & 1\,327  & \underline{73.46} & 1\,331  & \textbf{94.02} & 1\,332  \\
LoRA                & 368\,640 & \underline{81.15} & 2\,073  & 66.65 & 2\,123  & 81.25 & 2\,214  \\
\midrule
\multicolumn{8}{l}{\footnotesize \textit{Note:} Tip-Adapter-F param count scales as $16 \cdot C \cdot 512$; shown for EuroSAT ($C{=}10$).}\\
\multicolumn{8}{l}{\footnotesize \textit{Note:} LoRA is the only method modifying computation inside the backbone.}\\
\bottomrule
\end{tabular}
\end{table*}

\paragraph{Key observations.}
\begin{itemize}
  \item \textbf{Prompt ensembling is free.}  Averaging text features across seven templates adds $+$4.3 points on EuroSAT and $+$2.4 on DTD at zero parameter cost, establishing a stronger zero-shot baseline.

  \item \textbf{CoOp closes the gap efficiently.}  With only 8{,}192 parameters (0.0095\,\% of the backbone), CoOp achieves 75.87\,\% on EuroSAT, a $+$29.5 point improvement over zero-shot. On Flowers102, its 88.71\,\% surpasses even the linear probe.

  \item \textbf{CLIP-Adapter offers consistent gains.}  The vision adapter (131{,}712 params) outperforms CoOp on EuroSAT (+5.4) and DTD (+1.9), where the visual domain shift is largest.

  \item \textbf{Joint adaptation shows diminishing returns.} The CoOp\,+\,Adapter model performs comparably to the adapter alone, suggesting that both methods target overlapping subspaces of the classification gap.

  \item \textbf{Tip-Adapter-F dominates on accuracy.}  Fine-tuning the visual cache achieves the highest accuracy on EuroSAT (86.26\,\%) and Flowers102 (94.02\,\%), and the second-highest on DTD (73.46\,\%), at moderate parameter cost.

  \item \textbf{LoRA is competitive on EuroSAT but lags elsewhere.} Despite having the most parameters (368{,}640) and the highest memory footprint, LoRA does not generalise as well as lighter methods on DTD and Flowers102, likely due to overfitting on the 16-shot training set.

  \item \textbf{Optimal Transport fails.}  The Sinkhorn method scores \emph{below random chance} on Flowers102 (19.1\,\% for 102 classes). We analyse this failure in Section~\ref{sec:ot-analysis}.
\end{itemize}


\subsection{Accuracy vs.\ Trainable Parameters}

Figure~\ref{fig:acc-vs-params} plots accuracy as a function of trainable parameter count.  The curve shows a strong initial slope (zero-shot → CoOp) that flattens rapidly: the first 8{,}192 parameters buy the majority of the accuracy gain, and the marginal return per parameter decreases sharply thereafter.

\begin{figure}[t]
  \centering
  \includegraphics[width=\linewidth]{plots/accuracy_vs_params_unified.png}
  \caption{Top-1 accuracy vs.\ trainable parameter count across all methods
    and datasets.  The dashed line marks zero-shot performance.}
  \label{fig:acc-vs-params}
\end{figure}

\subsection{Training Cost and Resource Usage}

Table~\ref{tab:cost} summarises training time and peak GPU memory.
Tip-Adapter-F trains in under 150\,s (20 epochs on the cache only);
CoOp and CLIP-Adapter require 280--710\,s depending on the dataset; the
joint model's time is dominated by the CoOp-side gradient computation
through the text transformer.

\begin{table}[t]
\centering
\caption{Training time (seconds) for trainable methods.}
\label{tab:cost}
\small
\begin{tabular}{@{}lrrr@{}}
\toprule
\textbf{Method} & \textbf{EuroSAT} & \textbf{DTD} & \textbf{Flowers102} \\
\midrule
CoOp            & 357  & 285  & 705  \\
CLIP-Adapter    & 356  & 265  & 684  \\
CoOp+Adapter    & 420  & 289  & 713  \\
Tip-Adapter-F   &  44  &  99  & 146  \\
LoRA            & 370  & 290  & 689  \\
\bottomrule
\end{tabular}
\end{table}

\subsection{Analysis of Optimal Transport Failure}\label{sec:ot-analysis}

The Sinkhorn method under-performs zero-shot CLIP on every dataset.  We
identify three structural reasons:

\begin{enumerate}
  \item \textbf{Misaligned feature spaces.}  Patch tokens come from the
    penultimate layer of the vision transformer, \emph{before} contrastive
    alignment via the projection head.  Although we apply the output
    projection to all tokens, the contrastive loss only trained the pooled
    [CLS] token to be aligned with text; individual patch tokens inhabit
    a space that was never directly optimised for text matching.

  \item \textbf{Uninformative marginals.}  Uniform weighting treats
    background patches (sky, soil, water) with the same importance as
    object-relevant patches.  In satellite images, where the object of
    interest often occupies a small fraction of the $7 \times 7$ grid,
    the background dominates the transport plan.

  \item \textbf{Template-word dominance.}  The text-token set for each
    class shares template words (``a'', ``satellite'', ``image'', ``of''),
    which account for 5 of the $\sim$7 non-padding tokens.  The Sinkhorn
    distance is thus dominated by the cost of aligning image patches to
    these shared, non-discriminative tokens, drowning out the
    class-specific signal.
\end{enumerate}

These insights suggest concrete improvements (attention-weighted marginals,
class-name-only tokens) that we leave for future work.

\subsection{Ablation Studies}\label{sec:ablations}

\subsubsection{CoOp: Context Length and Learning Rate}

We sweep the context length $M \in \{4, 8, 16\}$ and learning rate
$\mathrm{lr} \in \{5 \times 10^{-4}, 10^{-3}, 2 \times 10^{-3},
5 \times 10^{-3}\}$.

% TODO: Insert ablation table from coop_results.json once you verify
% which numbers correspond to your final run.
% Example format:
% \begin{table}[t]
% \centering
% \caption{CoOp ablation: effect of context length $M$ and learning rate.}
% \label{tab:coop-ablation}
% \small
% \begin{tabular}{@{}lcccc@{}}
% \toprule
% M & lr=5e-4 & lr=1e-3 & lr=2e-3 & lr=5e-3 \\
% \midrule
% 4  & ... & ... & ... & ... \\
% 8  & ... & ... & ... & ... \\
% 16 & ... & ... & ... & ... \\
% \bottomrule
% \end{tabular}
% \end{table}

Longer context ($M{=}16$) consistently outperforms shorter ones,
suggesting that the additional capacity is not wasted in the 16-shot
regime.  The optimal learning rate is $2 \times 10^{-3}$ with SGD
and momentum, matching the CoOp paper's recommendation.

\subsubsection{CLIP-Adapter: Reduction Ratio and Alpha}

We sweep the reduction ratio $r \in \{2, 4, 8, 16\}$ and the blending
factor $\alpha \in \{0.1, 0.2, 0.3, 0.5, 0.7\}$.

% TODO: Insert ablation table from clip_adapter_results.json and
% clip_adapter_alpha_results.json once verified.
% \begin{table}[t]
% \centering
% \caption{CLIP-Adapter ablation: effect of reduction ratio $r$ and
%   blending factor $\alpha$.}
% \label{tab:adapter-ablation}
% \end{table}

$r{=}4$ (hidden dimension 128) and $\alpha{=}0.2$ form the best
combination.  Larger $\alpha$ values cause the adapter to overwrite
too much of the pretrained representation; smaller reduction ratios
introduce more parameters without a proportional accuracy gain.

\subsubsection{Unified vs.\ Class-Specific Context (CSC)}

Unified context ($M \times D$ shared parameters) and class-specific
context ($C \times M \times D$ separate parameters) yield comparable
accuracy in the 16-shot setting.  CSC has $C$ times more parameters
but does not improve on unified context, likely because the 16-shot
training set is too small to exploit the additional capacity.

\subsection{Qualitative Analysis}

Figure~\ref{fig:confusion} shows confusion matrices for the best-performing
methods on EuroSAT, highlighting that the most common failure mode is
confusion between ``annual crop land'' and ``permanent crop land''---two
visually similar categories that share agricultural patterns.

\begin{figure}[t]
  \centering
  \includegraphics[width=\linewidth]{plots/confusion_matrices.png}
  \caption{Confusion matrices for selected methods on EuroSAT.}
  \label{fig:confusion}
\end{figure}

Figure~\ref{fig:tsne} presents t-SNE visualisations of the adapted feature
spaces, showing that trainable methods produce more class-separable
embeddings than zero-shot CLIP.

\begin{figure}[t]
  \centering
  \includegraphics[width=\linewidth]{plots/tsne_features.png}
  \caption{t-SNE of image features: zero-shot vs.\ adapted methods.}
  \label{fig:tsne}
\end{figure}


% ============================================================================
% 5. CONCLUSION
% ============================================================================
\section{Conclusion}\label{sec:conclusion}

\subsection{Summary of Findings}

We have presented a unified, controlled comparison of seven
parameter-efficient adaptation strategies for CLIP, evaluated on three
visually diverse benchmarks with a common frozen ViT-B/32 backbone.
Our main findings are:

\begin{enumerate}
  \item \textbf{Extreme parameter efficiency is achievable.}
    CoOp's 8{,}192 learnable context tokens---less than 0.01\,\% of the
    backbone---recover a substantial fraction of the accuracy gap
    between zero-shot and full supervision.

  \item \textbf{Tip-Adapter-F offers the best accuracy--cost trade-off.}
    Fine-tuning a visual key--value cache achieves the highest accuracy on
    two of three datasets while training in a fraction of the time
    required by gradient-based methods.

  \item \textbf{Joint text--vision adaptation yields marginal gains.}
    Combining CoOp with CLIP-Adapter does not reliably outperform either
    component alone, suggesting that both methods exploit overlapping
    subspaces of the representation gap.

  \item \textbf{Token-level Optimal Transport fails without careful
    engineering.}  Naïve application of Sinkhorn distances to patch and
    word tokens produces accuracies below zero-shot, due to misaligned
    feature spaces, uninformative marginals, and template-word dominance.
\end{enumerate}

\subsection{Impact}

These results have practical implications for deploying CLIP-based
classifiers in resource-constrained settings.  For practitioners who need
domain-specific accuracy without the cost of full fine-tuning, our
analysis provides a clear ranking: \emph{start with prompt ensembling
(free), add CoOp if labels are available, and consider Tip-Adapter-F
when accuracy is paramount.}

\subsection{Limitations}

\begin{itemize}
  \item \textbf{No validation split.}  We use the test set for
    per-epoch model selection and hyperparameter tuning, following
    standard practice in the few-shot CLIP literature.  This means our
    reported accuracies are optimistic estimates.
  \item \textbf{LAION-2B weights, not OpenAI.}  We use the
    \texttt{laion2b\_s34b\_b79k} checkpoint rather than the original
    OpenAI weights.  Zero-shot baselines differ between the two.
  \item \textbf{Single backbone.}  All experiments use ViT-B/32;
    results may not transfer to larger architectures (ViT-L/14).
  \item \textbf{No data augmentation.}  Unlike the original CoOp and
    CLIP-Adapter papers, we do not apply random crop or horizontal flip
    during training.
  \item \textbf{Single seed.}  Most results use a single random seed.
    While training curves are stable, multi-seed confidence intervals
    would strengthen the conclusions.
\end{itemize}

\subsection{Future Work}

Several directions emerge naturally from our findings:

\begin{itemize}
  \item \textbf{Attention-weighted OT.}  Replace uniform marginals with
    attention-derived weights and restrict text tokens to class-specific
    words only, potentially rescuing the Optimal Transport approach.
  \item \textbf{CoCoOp.}  Extend CoOp with conditional context
    generation~\cite{zhou2022conditional} to improve cross-domain
    generalisation.
  \item \textbf{Larger backbones.}  Evaluate all methods on ViT-L/14 and
    ViT-H/14 to test whether the efficiency--accuracy trade-offs hold at
    scale.
  \item \textbf{Proper validation protocol.}  Introduce a held-out
    validation split to decouple model selection from test evaluation.
  \item \textbf{Data augmentation.}  Incorporate standard augmentations
    (random resized crop, horizontal flip) to match the original papers'
    protocols and improve few-shot accuracy.
\end{itemize}


% ============================================================================
% 6. CONTRIBUTIONS
% ============================================================================
\section{Contributions}\label{sec:contributions}

% TODO: Fill in specific contributions for each team member.
% The structure below follows the README; adjust as needed.

\begin{itemize}
  \item \textbf{Mattia:}
    \begin{itemize}
      \item Infrastructure: project architecture, \texttt{BaseCLIPWrapper}
        base class, shared evaluation pipeline (\texttt{engine.py}),
        dataset loaders, few-shot sampling protocol.
      \item Baselines: Zero-Shot CLIP, Zero-Shot with prompt ensembling,
        Linear Probe.
      \item Optimal Transport: Sinkhorn-based classification module,
        patch-token and text-token extraction, failure analysis.
      \item Paper writing: [TODO: specify sections written].
    \end{itemize}
  \item \textbf{Carlo:}
    \begin{itemize}
      \item CoOp: \texttt{PromptLearner} module, learnable text context,
        unified and class-specific variants.
      \item Joint model: \texttt{CoOpAdapterModel}, MRO resolution,
        optimizer group split.
      \item CoOp ablation experiments: context length, CSC, learning rate,
        shot sweeps.
      \item Paper writing: [TODO: specify sections written].
    \end{itemize}
  \item \textbf{Marco:}
    \begin{itemize}
      \item CLIP-Adapter: \texttt{VisionAdapterModule}, bottleneck MLP,
        residual blending with fixed and learnable $\alpha$.
      \item Tip-Adapter \& Tip-Adapter-F: key--value cache, cache
        fine-tuning.
      \item LoRA: \texttt{LinearLoRA} module, injection into the vision
        tower's MLP layers.
      \item CLIP-Adapter ablation experiments: reduction ratio sweep,
        $\alpha$ sweep.
      \item Paper writing: [TODO: specify sections written].
    \end{itemize}
\end{itemize}


% ============================================================================
% REFERENCES
% ============================================================================

{\small
\bibliographystyle{ieee_fullname}
\bibliography{references}

}

\end{document}
